% concatenated from Manuscript_DHKQ.tex
\documentclass[12pt,reqno]{amsart}
\usepackage[utf8]{inputenc}
\usepackage[margin=1.5 in]{geometry}
\usepackage[lite]{amsrefs}
\usepackage{amsfonts}
\usepackage{mathrsfs}
\usepackage{amscd}
%\usepackage{biblatex}
%\addbibresource{biblio.bib}

\usepackage{physics}
\usepackage{amsmath}
\usepackage{mathtools}
\usepackage{amssymb}
\usepackage{multicol}
\usepackage{graphicx}
\usepackage{wrapfig}
\usepackage{float}
\usepackage[english]{babel}
\usepackage[utf8]{inputenc}
\usepackage{fancyhdr}
\usepackage{amsthm}
\usepackage[all]{xy}
\usepackage{tikz-cd}
\usepackage{bbold}
\usepackage{hyperref}
\usepackage{tabularx,booktabs,caption,ragged2e}

\newcolumntype{L}{>{\RaggedRight\arraybackslash}X}
\captionsetup{skip=0.333\baselineskip}

\newtheorem{thm}{Theorem}[section]
\newtheorem{cor}[thm]{Corollary}
\newtheorem{lemma}[thm]{Lemma}
\newtheorem{propn}[thm]{Proposition}
\newtheorem{conj}[thm]{Conjecture}
\newtheorem*{question}{Question}

\theoremstyle{definition}
\newtheorem{defn}[thm]{Definition}
\newtheorem{example}[thm]{Example}

\theoremstyle{remark}
\newtheorem{remark}[thm]{Remark}
\newtheorem{claim}[thm]{Claim}

\newcommand{\g}{\mathfrak{g}}
\newcommand{\h}{\mathfrak{h}}
\newcommand{\p}{\mathfrak{p}}
\newcommand{\R}{\mathbb{R}}
\newcommand{\C}{\mathbb{C}}
\newcommand{\N}{\mathbb{N}}
\newcommand{\T}{\mathbb{T}}
\newcommand{\Z}{\mathbb{Z}}
\newcommand{\HS}{\mathcal{H}}
\newcommand{\B}{\mathcal{B}}
\newcommand{\A}{\mathcal{A}}
\newcommand{\E}{\mathcal{E}}
\newcommand{\Ad}{\operatorname{Ad}}
\newcommand{\ad}{\operatorname{ad}}
\newcommand{\GL}{\operatorname{GL}}
\newcommand{\SL}{\operatorname{SL}}
\newcommand{\Or}{\operatorname{O}}
\newcommand{\SO}{\operatorname{SO}}
\newcommand{\U}{\operatorname{U}}
\newcommand{\SU}{\operatorname{SU}}
\newcommand{\SP}{\operatorname{SP}}
\newcommand{\gl}{\mathfrak{gl}}
\newcommand{\sli}{\mathfrak{sl}}
\newcommand{\ort}{\mathfrak{o}}
\newcommand{\so}{\mathfrak{so}}
\newcommand{\uni}{\mathfrak{u}}
\newcommand{\su}{\mathfrak{su}}
\newcommand{\smp}{\mathfrak{sp}}
\newcommand{\n}{\nabla}

\newcommand{\PS}{\mathbb{P}}
\newcommand{\Sy}{\mathbb{S}}
\newcommand{\Ts}{\mathbb{T}}
\newcommand{\PSE}{\widehat{\mathbb{P}}}
\newcommand{\PW}{P_{\omega,t}}
\newcommand{\Rw}{R_{\omega,\lambda}}
\newcommand{\Aw}{A_{\omega,\lambda}}
\newcommand{\Sw}{S_\omega}
\newcommand{\RT}{R_{\omega,t}}
\newcommand{\TM}{\mathbb{T}M}
\newcommand{\smooth}{C_c^\infty(\R^n)}
\newcommand{\asmooth}{C_c^\infty(\R^n,\A)}
\newcommand{\esmooth}{C_c^\infty(\R^n,\E)}

\title{Finite-Rank Lie Algebroids for Singular Foliations of Prescribed Vanishing Order}

% \title{A Left-Symmetric Algebroid for Singular Foliations Vanishing to Prescribed Order}

\author{Austin Davis}
\address{Department of Mathematics, Penn State University}
\email{\href{mailto:austinhuntdavis@psu.edu}{austinhuntdavis@psu.edu}}

\author{Cheng-Yi Hung}
\address{Department of Mathematics, National Tsing Hua University}
\email{\href{mailto:chengyi@m113.nthu.edu.tw}{chengyi@m113.nthu.edu.tw}}

\author{Sangjun Ko}
\address{Department of Mathematics, Penn State University}
\email{\href{mailto:duck@psu.edu}{duck@psu.edu}}

\author{Kaichuan Qi}
\address{Department of Mathematics, Penn State University}
\email{\href{mailto:kaichuan@psu.edu}{kaichuan@psu.edu}}


% \thanks{Davis, Ko, and Qi were partially supported by the National Science Foundation
% (award DMS-2302447) and (award DMS-2201270).}

\begin{document}


\begin{abstract}
We construct a family of finite-rank Lie algebroids realizing, through
their anchor maps, the singular foliations of vector fields vanishing
to a prescribed order at the origin. The construction is obtained from
the sub-adjacent Lie algebroids of a class of left-symmetric
algebroids.
\end{abstract}

\maketitle
\section{Introduction}

It is a general open problem to determine which singular foliations
can be realized as the image sheaf of the anchor of a finite-rank Lie
algebroid. This problem was explicitly raised by
Androulidakis and Zambon \cite{AZ} and has since been highlighted in the
recent account of Laurent-Gengoux, Louis, and Ryvkin \cite{LGLR}. In this
note, we give a construction of finite-rank Lie algebroids
realizing a natural family of singular foliations on finite-dimensional
real vector spaces, consisting of vector fields with a prescribed order
of vanishing at the origin.

Let \(V=\mathbb{R}^n\), with standard coordinates
\(x_1,\ldots,x_n\), and fix an integer \(k\geq 1\).
Denote by \(C^\infty_V\) the sheaf of smooth functions on \(V\),
by \(\mathfrak X_V\) the sheaf of smooth vector fields on \(V\),
and by
\[
\mathcal I_0\subset C^\infty_V
\]
the ideal sheaf of the origin. Consider the singular foliation
\[
\mathcal F_k
=
\mathcal I_0^k\mathfrak X_V.
\]

For a multi-index
\(
\alpha=(\alpha_1,\ldots,\alpha_n)\in\mathbb{N}^n,
\)
write
\[
|\alpha|
=
\alpha_1+\cdots+\alpha_n,
\qquad
x^\alpha
=
x_1^{\alpha_1}\cdots x_n^{\alpha_n}.
\]
By the iterated Hadamard lemma, the ideal sheaf
\(\mathcal I_0^k\) is generated, as a \(C^\infty_V\)-module, by the
monomials
\[
x^\alpha,
\qquad
|\alpha|=k.
\]
Consequently, \(\mathcal F_k\) is globally generated, as a sheaf of
\(C^\infty_V\)-modules, by
\[
x^\alpha\frac{\partial}{\partial x_j},
\qquad
|\alpha|=k,
\qquad
1\leq j\leq n.
\]

\begin{question}
Is this singular foliation the image of the anchor of a
finite-rank Lie algebroid?
\end{question}

For \(k=1\), such a realization is provided by the standard action
Lie algebroid associated with the linear action of
\(\mathfrak{gl}(V)\) on \(V\). The cases \(n\geq 2\) and \(k\geq 2\)
are more subtle. In recent work, Louis \cite{Louis} considered the
same natural presentation of \(\mathcal F_k\) by a finite-rank
anchored vector bundle and constructed an almost Lie algebroid
structure on it, while noting that for \(k\neq 1\) it remained open whether $\mathcal{F}_k$ is the anchor image of a Lie algebroid.

We answer this question positively. More precisely, we construct a
family of finite-rank Lie algebroids whose anchor image is precisely
the singular foliation above. For the special case of \(n=k=2\), this
gives a positive answer to Question 1.8 of \cite{LGLR}. It also gives
a positive answer to the equivalent linear Poisson extension problem
raised by Lerman and Louis \cite{LL}.

Consider the trivial vector bundle
\[
A
=
V\times\bigl(S^kV^*\otimes V\bigr)
\longrightarrow V.
\]
Its rank is
\[
\operatorname{rank}(A)
=
n\binom{n+k-1}{k}.
\]
A section \(a\in\Gamma(A)\) is a smooth family
\(
x\longmapsto a_x,
\)
where
\[
a_x\colon V^k\longrightarrow V
\]
is a symmetric \(k\)-linear map. We first construct an anchor map.

Let \(\E\in\mathfrak X(V)\) be the Euler vector field. There is a natural
anchor
\[
\rho\colon A\longrightarrow TV
\]
defined on sections by
\[
\rho(a)
=
a(\underbrace{\E,\ldots,\E}_{k\text{ times}}),
\quad\text{equivalently,}\quad
\rho(a)_x=a_x(x,\ldots,x).
\]

For a multi-index \(\alpha\in\mathbb{N}^n\) with \(|\alpha|=k\), let
\[
dx^{\odot\alpha}
=
dx_1^{\odot\alpha_1}
\odot\cdots\odot
dx_n^{\odot\alpha_n}
\in S^kV^*
\]
denote the normalized symmetric tensor satisfying
\(
dx^{\odot\alpha}(x,\ldots,x)=x^\alpha.
\)
Then
\[
\rho\left(
dx^{\odot\alpha}\otimes
\frac{\partial}{\partial x_j}
\right)
=
x^\alpha\frac{\partial}{\partial x_j}.
\]
It follows that, for every open subset \(U\subseteq V\),
\[
\rho\bigl(\Gamma(U,A)\bigr)
=
\mathcal F_k(U).
\]

It remains to construct a Lie algebroid bracket on \(A\) compatible
with this anchor. We make use of the notion of left-symmetric
algebroids \cite{LSBC}, and construct a non-skew-symmetric product
\[
\triangleright\colon
\Gamma(A)\times\Gamma(A)
\longrightarrow\Gamma(A).
\]
We show that it defines a left-symmetric algebroid structure. The
associated sub-adjacent Lie algebroid has bracket
\[
[a,b]_A
=
a\triangleright b-b\triangleright a.
\]

The guiding requirement in the construction is the compatibility
condition
\[
\rho(a\triangleright b)
=
D_{\rho(a)}\rho(b),
\]
where \(D\) is the standard flat connection on \(V\). The
left-symmetric product is given by
\[
\begin{aligned}
(a\triangleright b)(u_1,\ldots,u_k)
={}&
(D_{\rho(a)}b)(u_1,\ldots,u_k)
\\
&+
\sum_{i=1}^k
b(u_1,\ldots,u_{i-1},
\Lambda_a u_i,
u_{i+1},\ldots,u_k),
\end{aligned}
\]
where
\[
\Lambda_a(u)
=
a(\underbrace{\E,\ldots,\E}_{k-1\text{ times}},u).
\]



\section{A class of left-symmetric algebroids}



 
\subsection{Construction of the left-symmetric algebroid}
We first recall basics of left-symmetric algebroids, following Liu, Sheng, Bai, and Chen \cite{LSBC}.
\begin{defn}
Let \(A\to M\) be an anchored vector bundle with anchor
\(\rho\colon A\to TM\). A \emph{left-symmetric algebroid structure}
on \(A\) is an \(\R\)-bilinear product
\[
\triangleright\colon
\Gamma(A)\times\Gamma(A)
\longrightarrow\Gamma(A)
\]
such that, for $f \in C^{\infty}(M)$ and $a,b,c \in \Gamma(A)$,
\begin{align}
(fa)\triangleright b
&=
f(a\triangleright b),
\label{eq:LSA-first-Leibniz}
\\
a\triangleright(fb)
&=
f(a\triangleright b)+\rho(a)[f]\,b,
\label{eq:LSA-second-Leibniz}
\end{align}
and whose associator
\begin{equation}
\operatorname{Ass}(a,b,c)
=
(a\triangleright b)\triangleright c
-
a\triangleright(b\triangleright c)
\label{eq:associator-definition}
\end{equation}
is symmetric in its first two arguments:
\begin{equation}
\operatorname{Ass}(a,b,c)
=
\operatorname{Ass}(b,a,c).
\label{eq:left-symmetric-identity}
\end{equation}
\end{defn}

\begin{propn}
Assume $(A,\rho,\triangleright)$ is a left-symmetric algebroid. Define
\[
[a,b]_A
=
a\triangleright b-b\triangleright a.
\]
Then
\[
\bigl(A,\rho,[\,\cdot,\cdot\,]_A\bigr)
\]
is a Lie algebroid.
\end{propn}



Now we construct the product $\triangleright$ for the left-symmetric algebroid. We regard \(V\) as an affine manifold equipped with its canonical
flat and torsion-free connection \(D\). We use the same notation
for the induced flat connections on all trivial tensor bundles over
\(V\).  We seek a lift satisfying
\[
\rho(a\triangleright b)
=
D_{\rho(a)}\rho(b).
\]
Since
\[
\rho(b)
=
b(\underbrace{\E,\ldots,\E}_{k\text{ times}}),
\]
by the Leibniz rule for the covariant derivative of tensor contractions,
%isn't this Leibniz?
\[
\begin{aligned}
D_{\rho(a)}\rho(b)
={}&
(D_{\rho(a)}b)(\E,\ldots,\E)
\\
&+
\sum_{i=1}^k
b(\E,\ldots,
\underbrace{D_{\rho(a)}\E}_{i\text{-th entry}},
\ldots,\E).
\end{aligned}
\]
The Euler vector field satisfies
\[
D_{\rho(a)}\E=\rho(a).
\]

For every \(a\in\Gamma(A)\), define an endomorphism field
\[
\Lambda_a
\in
C^\infty\bigl(V,\operatorname{End}(V)\bigr)
\]
by
\begin{equation}
\Lambda_a(u)
=
a(\underbrace{\E,\ldots,\E}_{k-1\text{ times}},u).
\label{eq:Lambda}
\end{equation}
Then
\begin{equation}\label{eq:Lambda-anchor}
\Lambda_a\E
=
\rho(a)
=
D_{\rho(a)}\E.    
\end{equation}

This suggests defining
\begin{equation}\label{eq:left-symmetric-product-expanded}
\begin{aligned}
&(a\triangleright b)(u_1,\ldots,u_k)
\\
={}&
(D_{\rho(a)}b)(u_1,\ldots,u_k)
\\
&+
\sum_{i=1}^k
b(u_1,\ldots,u_{i-1},
\Lambda_a u_i,
u_{i+1},\ldots,u_k).
\end{aligned}    
\end{equation}

By construction,
\begin{equation}
\rho(a\triangleright b)
=
D_{\rho(a)}\rho(b),
\label{eq:rho-product-identity}
\end{equation}

% -----

% We regard \(V\) as an affine manifold
% equipped with its canonical flat and torsion-free connection
% \[
% D\colon
% \mathfrak X(V)\times\mathfrak X(V)
% \longrightarrow
% \mathfrak X(V).
% \]
% We use the same notation \(D\) for the induced flat connections
% on all trivial tensor bundles over \(V\).

% Let
% \[
% \E\in\mathfrak X(V)
% \]
% be the Euler vector field,
% \[
% \E_x=x
% \qquad
% \text{for every }x\in V.
% \]
% Since the map \(x\mapsto x\) is the identity map of \(V\), one has
% \begin{equation}
% D_\xi\E=\xi
% \label{eq:Euler-derivative}
% \end{equation}
% for every vector field \(\xi\in\mathfrak X(V)\).

% Consider the trivial vector bundle
% \begin{equation}
% A
% =
% V\times\bigl(S^kV^*\otimes V\bigr)
% \longrightarrow V.
% \label{eq:bundle-A}
% \end{equation}


% Define an anchor
% \[
% \rho\colon A\longrightarrow TV
% \]
% by
% \begin{equation}
% \rho(a)
% =
% a(\underbrace{\E,\ldots,\E}_{k\text{ times}}).
% \label{eq:anchor}
% \end{equation}
% Equivalently,
% \[
% \rho(a)_x=a_x(x,\ldots,x).
% \]


% Thus
% \[
% (\Lambda_a)_x(u)
% =
% a_x(\underbrace{x,\ldots,x}_{k-1\text{ times}},u).
% \]
% In particular,
% \begin{equation}
% \Lambda_a(\E)=\rho(a).
% \label{eq:Lambda-anchor}
% \end{equation}

% For an endomorphism field
% \[
% \Phi
% \in
% C^\infty\bigl(V,\operatorname{End}(V)\bigr),
% \]
% define
% \[
% \mathcal P_\Phi\colon
% \Gamma(A)\longrightarrow\Gamma(A)
% \]
% by
% \begin{equation}
% \begin{aligned}
% &(\mathcal P_\Phi b)(u_1,\ldots,u_k)
% \\
% &\qquad=
% \sum_{i=1}^k
% b(u_1,\ldots,u_{i-1},
% \Phi u_i,
% u_{i+1},\ldots,u_k).
% \end{aligned}
% \label{eq:P-Phi}
% \end{equation}

% We define a non-skew-symmetric product
% \[
% \triangleright\colon
% \Gamma(A)\times\Gamma(A)
% \longrightarrow\Gamma(A)
% \]
% by
% \begin{equation}
% a\triangleright b
% =
% D_{\rho(a)}b+\mathcal P_{\Lambda_a}b.
% \label{eq:left-symmetric-product}
% \end{equation}
% Equivalently,
% \begin{equation}
% \begin{aligned}
% &(a\triangleright b)(u_1,\ldots,u_k)
% \\
% ={}&
% (D_{\rho(a)}b)(u_1,\ldots,u_k)
% \\
% &+
% \sum_{i=1}^k
% b(u_1,\ldots,u_{i-1},
% \Lambda_a u_i,
% u_{i+1},\ldots,u_k).
% \end{aligned}
% \label{eq:left-symmetric-product-expanded}
% \end{equation}

We also write
\begin{equation}
\nabla_ab:=a\triangleright b.
\label{eq:nabla-product}
\end{equation}

For an endomorphism field
\[
\Phi\in C^\infty\bigl(V,\operatorname{End}(V)\bigr),
\]
define
\[
\mathcal P_\Phi\colon\Gamma(A)\longrightarrow\Gamma(A)
\]
by
\begin{equation}
(\mathcal P_\Phi b)(u_1,\ldots,u_k)
=
\sum_{i=1}^k
b(u_1,\ldots,u_{i-1},
\Phi u_i,
u_{i+1},\ldots,u_k).
\label{eq:P-Phi}
\end{equation}
With this notation,
\[
a\triangleright b
=
D_{\rho(a)}b+\mathcal P_{\Lambda_a}b.
\]

The main result is
\begin{thm}\label{thm:Main}
The product
\[
a\triangleright b
=
D_{\rho(a)}b+\mathcal P_{\Lambda_a}b
\]
defines a left-symmetric algebroid structure on the anchored vector
bundle \(A\to V\).
\end{thm}



\subsection{A Leibniz-rule identity, and commutators of the tensorial operators}

We prepare the proof of Theorem \ref{thm:Main} with the following 2 lemmas.

\begin{lemma}
For all \(a,b\in\Gamma(A)\), we have
\begin{equation}
\Lambda_{a\triangleright b}
=
D_{\rho(a)}\Lambda_b
+
\Lambda_b\circ\Lambda_a.
\label{eq:Lambda-product-identity}
\end{equation}
\end{lemma}

\begin{proof}
% By definition,
% \[
% \rho(a\triangleright b)
% =
% (a\triangleright b)
% (\underbrace{\E,\ldots,\E}_{k\text{ times}}).
% \]
% Using \eqref{eq:left-symmetric-product-expanded}, one obtains
% \[
% \begin{aligned}
% \rho(a\triangleright b)
% ={}&
% (D_{\rho(a)}b)
% (\underbrace{\E,\ldots,\E}_{k\text{ times}})
% \\
% &+
% \sum_{i=1}^k
% b(\E,\ldots,
% \underbrace{\Lambda_a\E}_{i\text{-th entry}},
% \ldots,\E).
% \end{aligned}
% \]
% By \eqref{eq:Lambda-anchor},
% \[
% \Lambda_a\E=\rho(a).
% \]
% Therefore, using the symmetry of \(b\),
% \begin{equation}
% \begin{aligned}
% \rho(a\triangleright b)
% ={}&
% (D_{\rho(a)}b)
% (\underbrace{\E,\ldots,\E}_{k\text{ times}})
% \\
% &+
% k\,b\bigl(
% \rho(a),
% \underbrace{\E,\ldots,\E}_{k-1\text{ times}}
% \bigr).
% \end{aligned}
% \label{eq:rho-product-first}
% \end{equation}

% On the other hand,
% \[
% \rho(b)
% =
% b(\underbrace{\E,\ldots,\E}_{k\text{ times}}).
% \]
% Using \(D_{\rho(a)}\E=\rho(a)\), the chain rule gives
% \[
% \begin{aligned}
% D_{\rho(a)}\rho(b)
% ={}&
% D_{\rho(a)}
% \bigl(
% b(\underbrace{\E,\ldots,\E}_{k\text{ times}})
% \bigr)
% \\
% ={}&
% (D_{\rho(a)}b)
% (\underbrace{\E,\ldots,\E}_{k\text{ times}})
% \\
% &+
% \sum_{i=1}^k
% b(\E,\ldots,
% \underbrace{D_{\rho(a)}\E}_{i\text{-th entry}},
% \ldots,\E)
% \\
% ={}&
% (D_{\rho(a)}b)
% (\underbrace{\E,\ldots,\E}_{k\text{ times}})
% \\
% &+
% k\,b\bigl(
% \rho(a),
% \underbrace{\E,\ldots,\E}_{k-1\text{ times}}
% \bigr).
% \end{aligned}
% \]
% Comparing this with \eqref{eq:rho-product-first} proves
% \eqref{eq:rho-product-identity}.
 Let \(u\in V\)
be regarded as a constant vector field. Then
\[
\begin{aligned}
\Lambda_{a\triangleright b}(u)
={}&
(a\triangleright b)
(\underbrace{\E,\ldots,\E}_{k-1\text{ times}},u)
\\
={}&
(D_{\rho(a)}b)
(\underbrace{\E,\ldots,\E}_{k-1\text{ times}},u)
\\
&+
\sum_{i=1}^{k-1}
b(\E,\ldots,
\underbrace{\Lambda_a\E}_{i\text{-th entry}},
\ldots,\E,u)
\\
&+
b(\underbrace{\E,\ldots,\E}_{k-1\text{ times}},
\Lambda_a u).
\end{aligned}
\]
Using \(\Lambda_a\E=\rho(a)\) and the symmetry of \(b\), this becomes
\begin{equation}
\begin{aligned}
\Lambda_{a\triangleright b}(u)
={}&
(D_{\rho(a)}b)
(\underbrace{\E,\ldots,\E}_{k-1\text{ times}},u)
\\
&+
(k-1)b\bigl(
\rho(a),
\underbrace{\E,\ldots,\E}_{k-2\text{ times}},
u
\bigr)
\\
&+
\Lambda_b(\Lambda_a u).
\end{aligned}
\label{eq:Lambda-product-first}
\end{equation}
When \(k=1\), the middle term is absent.

Since \(u\) is constant,
\[
D_{\rho(a)}u=0.
\]
Therefore
\[
\begin{aligned}
\bigl(D_{\rho(a)}\Lambda_b\bigr)(u)
={}&
D_{\rho(a)}\bigl(\Lambda_b(u)\bigr)
\\
={}&
D_{\rho(a)}
\bigl(
b(\underbrace{\E,\ldots,\E}_{k-1\text{ times}},u)
\bigr)
\\
={}&
(D_{\rho(a)}b)
(\underbrace{\E,\ldots,\E}_{k-1\text{ times}},u)
\\
&+
(k-1)b\bigl(
\rho(a),
\underbrace{\E,\ldots,\E}_{k-2\text{ times}},
u
\bigr).
\end{aligned}
\]
Again, the last term is absent when \(k=1\).
Substituting this into \eqref{eq:Lambda-product-first} yields
\[
\Lambda_{a\triangleright b}(u)
=
\bigl(D_{\rho(a)}\Lambda_b\bigr)(u)
+
(\Lambda_b\circ\Lambda_a)(u).
\]
Since this holds for every constant vector \(u\), we obtain
\eqref{eq:Lambda-product-identity}.
\end{proof}



For operators on \(\Gamma(A)\), we use the convention
\[
[\mathcal D_1,\mathcal D_2]
=
\mathcal D_1\circ\mathcal D_2
-
\mathcal D_2\circ\mathcal D_1.
\]
For endomorphism fields, composition is written in the usual order:
\[
(\Phi\circ\Psi)(u)=\Phi(\Psi(u)).
\]

\begin{lemma}
For all \(a,b\in\Gamma(A)\),
\begin{align}
[D_{\rho(a)},D_{\rho(b)}]
&=
D_{[\rho(a),\rho(b)]},
\label{eq:DD-commutator}
\\
[D_{\rho(a)},\mathcal P_{\Lambda_b}]
&=
\mathcal P_{D_{\rho(a)}\Lambda_b},
\label{eq:DP-commutator}
\\
[\mathcal P_{\Lambda_a},\mathcal P_{\Lambda_b}]
&=
\mathcal P_{
\Lambda_b\circ\Lambda_a
-
\Lambda_a\circ\Lambda_b
}.
\label{eq:PP-commutator}
\end{align}
\end{lemma}

\begin{proof}
Identity \eqref{eq:DD-commutator} follows from the flatness of \(D\).

Let \(c\in\Gamma(A)\), and let
\(u_1,\ldots,u_k\in V\) be constant vector fields. By definition,
\[
(\mathcal P_{\Lambda_b}c)(u_1,\ldots,u_k)
=
\sum_{i=1}^k
c(u_1,\ldots,u_{i-1},
\Lambda_bu_i,
u_{i+1},\ldots,u_k).
\]
Differentiating along \(\rho(a)\), one obtains
\[
\begin{aligned}
&D_{\rho(a)}
\bigl(\mathcal P_{\Lambda_b}c\bigr)
(u_1,\ldots,u_k)
\\
={}&
\sum_{i=1}^k
(D_{\rho(a)}c)
(u_1,\ldots,u_{i-1},
\Lambda_bu_i,
u_{i+1},\ldots,u_k)
\\
&+
\sum_{i=1}^k
c(u_1,\ldots,u_{i-1},
(D_{\rho(a)}\Lambda_b)u_i,
u_{i+1},\ldots,u_k).
\end{aligned}
\]
Meanwhile,
\[
\begin{aligned}
&\mathcal P_{\Lambda_b}
\bigl(D_{\rho(a)}c\bigr)
(u_1,\ldots,u_k)
\\
={}&
\sum_{i=1}^k
(D_{\rho(a)}c)
(u_1,\ldots,u_{i-1},
\Lambda_bu_i,
u_{i+1},\ldots,u_k).
\end{aligned}
\]
Subtracting gives
\[
\begin{aligned}
&[D_{\rho(a)},\mathcal P_{\Lambda_b}]c
(u_1,\ldots,u_k)
\\
={}&
\sum_{i=1}^k
c(u_1,\ldots,u_{i-1},
(D_{\rho(a)}\Lambda_b)u_i,
u_{i+1},\ldots,u_k)
\\
={}&
\mathcal P_{D_{\rho(a)}\Lambda_b}c
(u_1,\ldots,u_k),
\end{aligned}
\]
which proves \eqref{eq:DP-commutator}.

For the last identity, expand
\[
\mathcal P_{\Lambda_a}
\mathcal P_{\Lambda_b}c
\]
and
\[
\mathcal P_{\Lambda_b}
\mathcal P_{\Lambda_a}c.
\]
The terms in which \(\Lambda_a\) and \(\Lambda_b\) act in different
tensor slots occur identically in the two expansions and cancel
after subtraction. The remaining terms are those in which the two
endomorphisms act in the same tensor slot. Hence
\[
\begin{aligned}
&[\mathcal P_{\Lambda_a},\mathcal P_{\Lambda_b}]c
(u_1,\ldots,u_k)
\\
={}&
\sum_{i=1}^k
c\bigl(
u_1,\ldots,u_{i-1},
(\Lambda_b\circ\Lambda_a
-
\Lambda_a\circ\Lambda_b)u_i,
u_{i+1},\ldots,u_k
\bigr)
\\
={}&
\mathcal P_{
\Lambda_b\circ\Lambda_a
-
\Lambda_a\circ\Lambda_b
}c(u_1,\ldots,u_k).
\end{aligned}
\]
This proves \eqref{eq:PP-commutator}.
\end{proof}

\subsection{Proof of Theorem \ref{thm:Main}}
\begin{proof}
We first verify the two Leibniz identities.

Let \(f\in C^\infty(V)\). Since
\[
\rho(fa)=f\rho(a),
\qquad
\Lambda_{fa}=f\Lambda_a,
\]
one has
\[
\begin{aligned}
(fa)\triangleright b
&=
D_{\rho(fa)}b+\mathcal P_{\Lambda_{fa}}b
\\
&=
D_{f\rho(a)}b+\mathcal P_{f\Lambda_a}b
\\
&=
fD_{\rho(a)}b+f\mathcal P_{\Lambda_a}b
\\
&=
f(a\triangleright b).
\end{aligned}
\]
Also,
\[
\begin{aligned}
a\triangleright(fb)
&=
D_{\rho(a)}(fb)+\mathcal P_{\Lambda_a}(fb)
\\
&=
\rho(a)[f]\,b
+
fD_{\rho(a)}b
+
f\mathcal P_{\Lambda_a}b
\\
&=
f(a\triangleright b)+\rho(a)[f]\,b.
\end{aligned}
\]
Thus \eqref{eq:LSA-first-Leibniz} and
\eqref{eq:LSA-second-Leibniz} hold.

Define the commutator
\begin{equation}
[a,b]_A
=
a\triangleright b-b\triangleright a.
\label{eq:subadjacent-bracket}
\end{equation}
Using \eqref{eq:rho-product-identity}, one obtains
\[
\begin{aligned}
\rho([a,b]_A)
&=
\rho(a\triangleright b)
-
\rho(b\triangleright a)
\\
&=
D_{\rho(a)}\rho(b)
-
D_{\rho(b)}\rho(a).
\end{aligned}
\]
Since \(D\) is torsion-free,
\[
D_{\rho(a)}\rho(b)
-
D_{\rho(b)}\rho(a)
=
[\rho(a),\rho(b)].
\]
Therefore
\begin{equation}
\rho([a,b]_A)
=
[\rho(a),\rho(b)].
\label{eq:anchor-bracket-compatibility}
\end{equation}

We next establish the flatness identity
\begin{equation}
[\nabla_a,\nabla_b]
=
\nabla_{[a,b]_A}.
\label{eq:flatness-identity}
\end{equation}
Recall that, as an operator on \(\Gamma(A)\), 
\[
\nabla_a
=
D_{\rho(a)}+\mathcal P_{\Lambda_a}.
\]
Therefore
\[
\begin{aligned}
[\nabla_a,\nabla_b]
={}&
[D_{\rho(a)}+\mathcal P_{\Lambda_a},
D_{\rho(b)}+\mathcal P_{\Lambda_b}]
\\
={}&
[D_{\rho(a)},D_{\rho(b)}]
+
[D_{\rho(a)},\mathcal P_{\Lambda_b}]
\\
&+
[\mathcal P_{\Lambda_a},D_{\rho(b)}]
+
[\mathcal P_{\Lambda_a},\mathcal P_{\Lambda_b}].
\end{aligned}
\]
Using
\[
[\mathcal P_{\Lambda_a},D_{\rho(b)}]
=
-[D_{\rho(b)},\mathcal P_{\Lambda_a}],
\]
together with
\eqref{eq:DD-commutator}--\eqref{eq:PP-commutator}, gives
\begin{align}
[\nabla_a,\nabla_b]
={}&
D_{[\rho(a),\rho(b)]}
\nonumber\\
&+
\mathcal P_{
D_{\rho(a)}\Lambda_b
-
D_{\rho(b)}\Lambda_a
+
\Lambda_b\circ\Lambda_a
-
\Lambda_a\circ\Lambda_b
}.
\label{eq:operator-commutator-expanded}
\end{align}

By \eqref{eq:Lambda-product-identity},
\[
\Lambda_{a\triangleright b}
=
D_{\rho(a)}\Lambda_b+\Lambda_b\circ\Lambda_a,
\]
and hence
\begin{align}
\Lambda_{[a,b]_A}
={}&
D_{\rho(a)}\Lambda_b
-
D_{\rho(b)}\Lambda_a
\nonumber\\
&+
\Lambda_b\circ\Lambda_a
-
\Lambda_a\circ\Lambda_b.
\label{eq:Lambda-bracket}
\end{align}
Combining \eqref{eq:anchor-bracket-compatibility},
\eqref{eq:operator-commutator-expanded}, and
\eqref{eq:Lambda-bracket}, we obtain
\[
\begin{aligned}
[\nabla_a,\nabla_b]
&=
D_{\rho([a,b]_A)}
+
\mathcal P_{\Lambda_{[a,b]_A}}
\\
&=
\nabla_{[a,b]_A}.
\end{aligned}
\]
This proves \eqref{eq:flatness-identity}.

It remains to verify the left-symmetric identity. By definition,
\[
\begin{aligned}
\operatorname{Ass}(a,b,c)
&=
(a\triangleright b)\triangleright c
-
a\triangleright(b\triangleright c)
\\
&=
\nabla_{\nabla_ab}c-\nabla_a\nabla_bc.
\end{aligned}
\]
Therefore
\[
\begin{aligned}
&\operatorname{Ass}(a,b,c)
-
\operatorname{Ass}(b,a,c)
\\
={}&
\nabla_{\nabla_ab}c
-
\nabla_a\nabla_bc
-
\nabla_{\nabla_ba}c
+
\nabla_b\nabla_ac
\\
={}&
\nabla_{\nabla_ab-\nabla_ba}c
-
[\nabla_a,\nabla_b]c
\\
={}&
\nabla_{[a,b]_A}c
-
[\nabla_a,\nabla_b]c.
\end{aligned}
\]
By the flatness identity \eqref{eq:flatness-identity}, the last
expression vanishes. Thus
\[
\operatorname{Ass}(a,b,c)
=
\operatorname{Ass}(b,a,c).
\]
Hence \((A,\triangleright,\rho)\) is a left-symmetric algebroid. This finishes the proof of Theorem \ref{thm:Main}.
\end{proof}

\subsection{The anchor image}

We now return to the singular foliation
\[
\mathcal F_k
=
\mathcal I_0^k\mathfrak X_V.
\]

\begin{propn}
The image of the anchor agrees with \(\mathcal F_k\). More precisely, for
every open subset \(U\subseteq V\),
\begin{equation}
\rho\bigl(\Gamma(U,A)\bigr)
=
\mathcal F_k(U),
\label{eq:sheaf-anchor-image}
\end{equation}
and
\begin{equation}
\rho\bigl(\Gamma_c(A)\bigr)
=
\mathcal F_k(V)\cap\mathfrak X_c(V).
\label{eq:compact-anchor-image}
\end{equation}
\end{propn}

\begin{proof}
The constant sections
\[
dx^{\odot\alpha}\otimes \frac{\partial}{\partial x_j},
\qquad
|\alpha|=k,\quad 1\leq j\leq n,
\]
form a global frame of \(A\), and their anchor images are
\[
x^\alpha\frac{\partial}{\partial x_j}.
\]
These are precisely the global generators of
\(\mathcal I_0^k\mathfrak X_V\) as a sheaf of
\(C^\infty_V\)-modules.

We now verify that this gives equality on sections over every
open subset \(U\subseteq V\). The inclusion
\[
\rho\bigl(\Gamma(U,A)\bigr)\subseteq \mathcal F_k(U)
\]
is immediate. Conversely, let \(X\in\mathcal F_k(U)\). Since the vector
fields
\[
x^\alpha\frac{\partial}{\partial x_j},
\qquad
|\alpha|=k,\quad 1\leq j\leq n,
\]
generate \(\mathcal F_k\) as a sheaf, there is an open cover
\(\{U_\ell\}\) of \(U\) and sections
\(a_\ell\in\Gamma(U_\ell,A)\) such that
\[
\rho(a_\ell)=X|_{U_\ell}.
\]
Choose a locally finite partition of unity
\(\{\chi_\ell\}\) subordinate to \(\{U_\ell\}\). Then
\[
a=\sum_\ell \chi_\ell a_\ell\in\Gamma(U,A)
\]
is well defined and, since the anchor is \(C^\infty\)-linear,
\[
\rho(a)
=
\sum_\ell \chi_\ell\rho(a_\ell)
=
\sum_\ell \chi_\ell X
=
X.
\]
Hence, for every open subset \(U\subseteq V\),
\[
\rho\bigl(\Gamma(U,A)\bigr)=\mathcal F_k(U).
\]

The compactly supported statement follows by the usual cutoff argument.
Indeed, if
\[
X=\rho(a)\in\mathfrak X_c(V)
\]
for some \(a\in\Gamma(A)\), choose
\(\chi\in C^\infty_c(V)\) equal to \(1\) on a neighborhood of
\(\operatorname{supp}(X)\). Then
\[
\chi a\in\Gamma_c(A),
\qquad
\rho(\chi a)=\chi X=X.
\]
Thus
\[
\rho\bigl(\Gamma_c(A)\bigr)
=
\rho\bigl(\Gamma(A)\bigr)\cap\mathfrak X_c(V)
=
\mathcal F_k(V)\cap\mathfrak X_c(V).
\]
\end{proof}


\begin{cor}
For every \(n\geq1\) and \(k\geq1\), the singular foliation
\[
\mathcal F_k
=
\mathcal I_0^k\mathfrak X_{\R^n}
\]
is the anchor image of the finite-rank Lie algebroid
\[
A
=
\R^n\times
\bigl(S^k(\R^n)^*\otimes\R^n\bigr)
\longrightarrow\R^n,
\]
whose rank is
\[
n\binom{n+k-1}{k}.
\]
\end{cor}

\begin{remark}
The problem considered in this note is also related to recent work
of Lavau \cite{L} and of Lerman--Louis \cite{LL}. We hope that the construction presented
here provides new examples, and offers a complementary perspective, for their theories.
\end{remark}    


\vspace{20pt}  

\noindent \textbf{Acknowledgments.} We thank Leonid Ryvkin for bringing this problem to our attention
during GAP 2026 at the Research Institute for Mathematical Sciences
(RIMS), Kyoto University. We are grateful to Ping Xu, Mathieu Sti\'{e}non, and Hsuan-Yi Liao  for fruitful discussions and useful comments.

% \vspace{20pt} 

% \noindent \textbf{Data Availability Statement.} 
% No data were used for the preceding work, as it is a mathematical result. Relevant materials may be found within the works cited.
% \vspace{20pt}  

 

% \noindent \textbf{Conflict of interest statement.} 
% On behalf of all authors, the corresponding author states that there is no conflict of interest. 



%\noindent \textbf{Data Availability Statement.} 
%No data was used for the research described in the article.

\begin{bibdiv}
\begin{biblist}

%\parskip = 0.9em %这里调整的是参考文献段落间的距离

\bib{AZ}{article}{
   author={Androulidakis, Iakovos},
   author={Zambon, Marco},
   title={Smoothness of holonomy covers for singular foliations and
   essential isotropy},
   journal={Math. Z.},
   volume={275},
   date={2013},
   number={3-4},
   pages={921--951},
   issn={0025-5874},
}



\bib{LGLR}{article}{
   author={Laurent-Gengoux, Camille},
   author={Louis, Ruben},
   author={Ryvkin, Leonid},
   title={State of the art and open questions},
   conference={
      title={Advances in Poisson geometry},
   },
   book={
      series={Adv. Courses Math. CRM Barcelona},
      publisher={Birkh\"auser/Springer, Cham},
   },
   isbn={978-3-031-86388-2},
   isbn={978-3-031-86657-9},
   date={[2025] \copyright 2025},
   pages={391--415},
}


\bib{L}{article}{
   author={Lavau, Sylvain},
   title={The modular class of a singular foliation},
   journal={J. Geom. Phys.},
   volume={192},
   date={2023},
   pages={Paper No. 104902, 32},
   issn={0393-0440},
}

\bib{LL}{article}{
  author={Lerman, Eugene},
  author={Louis, Ruben},
  title={Lie-Rinehart and Poisson algebras over {$C^{\infty}$}-rings}, 
	eprint={https://arxiv.org/pdf/2606.01388},
   % doi={10.48550/arXiv.2601.05183},
}

\bib{LSBC}{article}{
   author={Liu, Jiefeng},
   author={Sheng, Yunhe},
   author={Bai, Chengming},
   author={Chen, Zhiqi},
   title={Left-symmetric algebroids},
   journal={Math. Nachr.},
   volume={289},
   date={2016},
   number={14-15},
   pages={1893--1908},
   issn={0025-584X},
}

\bib{Louis}{article}{
   author={Louis, Ruben},
   title={On Nash resolution of (singular) Lie algebroids},
   journal={Math. Z.},
   volume={311},
   date={2025},
   number={2},
   pages={Paper No. 24, 27},
}






\end{biblist}
\end{bibdiv}

\end{document}


